\chapter{Langevin and Brownian dynamics}

\section{Validity of the Langevin approach}

As suggested, let’s consider a time scale of $\tau_{\mathrm{coll}} =
\qty{e-12}{s}$ for the heat bath fluctuations. We want our algorithm to have a
time step $\Delta t$ at least \num{10} times larger than $\tau_{\mathrm{coll}}$,
i.e. $\Delta t > \qty{e-11}{s}$. We also assume the characteristic time of the
dynamics $\tau$ to be of the order of $\num{100}\, \Delta t$.

$\tau$ can be related to the mass, characteristic length $\sigma$ and
characteristic energy -- which we take equal to the thermal energy
$k_{\mathrm{B}} T$ at $T = \qty{300}{K}$ -- of the particle as follows:
\begin{equation}
  \tau = \sqrt{\frac{m \sigma^{2}}{k_{\mathrm{B}} T}}.
\end{equation}
If we take a spherical particle with roughly the same density as water, we get
\begin{equation}
  \tau = \sqrt{\frac{\pi \rho \sigma^{5}}{6 \smash[b]{k_{\mathrm{B}}} T}}
  \approx (\qty{3.6e11}{s\per\meter^{5/2}}) \, \sigma^{5/2},
\end{equation}
which implies the inequality
\begin{equation}
  \sigma \approx (\qty{2.4e-5}{m\per\second^{2/5}})\, \tau^{2/5} >
  \qty{6e-9}{m}
  \label{eq:langevin-req}
\end{equation}
for the characteristic length scale of the particle.

For the second question, we start from the mean square displacement in a
diffusive regime,
\begin{equation}
  \expval{(x(t) - x_{0})^{2}} \approx 2Dt.
\end{equation}
To have a displacement of $\sigma$ in a time $\Delta t$ we need
\begin{equation}
  \sigma^{2} \approx 2 D \Delta t \implies
  D \approx \frac{\sigma^{2}}{2 \Delta t}.
\end{equation}
From the Stokes--Einstein relation,
\begin{equation}
  D = \frac{k_{\mathrm{B}} T}{3 \pi \eta \sigma},
\end{equation}
and taking $\eta = \qty{0.853e-3}{Pa.s}$ as the viscosity of water at
\qty{300}{K}, we get:
\begin{equation}
  \frac{k_{\mathrm{B}} T}{3 \pi \eta \sigma} \approx
  \frac{\sigma^{2}}{2 \Delta t} \implies \Delta t \approx
  \frac{3 \pi \eta \sigma^{3}}{2k_{\mathrm{B}} T} \approx
  (\qty{9.7e-17}{s\per m^{3}}) \, \sigma^{3}.
  \label{eq:timescale-brownian}
\end{equation}
With $\sigma = \qty{e-8}{m}$ we have a time scale $\Delta t \approx
\qty{e-6}{s}$, perfectly compatible with the requirement for the Langevin
description.

Lastly, the exercise asked to show that Brownian motion becomes negligible for
particles sizes above
$\sigma = \qty{5}{\micro m}$. If we assume the Brownian motion description is
valid up to a time scale on the order of seconds, and we set $\Delta t =
\qty{1}{s}$ in \cref{eq:timescale-brownian}, we get:
\begin{equation}
  \sigma = \biggl(\frac{2k_{\mathrm{B}} T \Delta t}{3 \pi \eta}\biggr)^{1/3}
  \approx \qty{1}{\micro m}.
\end{equation}
Beyond this scale, the Langevin description is no longer applicable, and
inertial forces become dominant.
